\documentclass[11pt,a4paper]{article}

% ── Compilar con: tectonic x5_alt_v2.tex  (motor XeTeX, UTF-8 nativo) ──
\usepackage{fontspec}
\usepackage{amsmath,amssymb}
\usepackage[margin=2.3cm]{geometry}
\usepackage{xcolor}
\usepackage{booktabs}
\usepackage{tabularx}
\usepackage{longtable}
\usepackage{listings}
\usepackage{enumitem}
\usepackage{underscore}
\usepackage{titlesec}
\usepackage[hidelinks]{hyperref}

% ── Colores (mismos que docs/plans/algoritmos.tex y formulas.tex) ──
\definecolor{cazul}{HTML}{2C5F8A}
\definecolor{cverde}{HTML}{2E8B57}
\definecolor{cnaranja}{HTML}{C1660A}
\definecolor{cmorado}{HTML}{6B4FA0}
\definecolor{cgris}{HTML}{5B6770}
\definecolor{cfondo}{HTML}{EDF2F7}
\definecolor{ccodefondo}{HTML}{F6F8FA}

\setcounter{secnumdepth}{3}
\setcounter{tocdepth}{2}

\titleformat{\paragraph}[hang]{\normalfont\normalsize\bfseries\color{cazul}}{}{0pt}{}
\titlespacing*{\paragraph}{0pt}{2.0ex plus 1ex minus .2ex}{0.8ex plus .2ex}

\renewcommand{\contentsname}{Índice}
\renewcommand{\tablename}{Tabla}

\lstset{
  basicstyle=\ttfamily\footnotesize,
  backgroundcolor=\color{ccodefondo},
  frame=single,
  rulecolor=\color{cgris!40},
  breaklines=true,
  columns=fullflexible,
  keepspaces=true,
  showstringspaces=false,
  xleftmargin=6pt,
  framexleftmargin=4pt,
}
\newcommand{\code}[1]{\texttt{\small #1}}

% ── Caja de nota (naranja) ──
\newenvironment{nota}{%
  \par\medskip\noindent
  \begin{tabular}{@{}p{0.4em}@{\hspace{8pt}}p{\dimexpr\linewidth-2em\relax}@{}}
  \textcolor{cnaranja}{\rule{3pt}{1.1em}} &
}{\end{tabular}\par\medskip}

% ── Caja de clave (morada) — equivalente a .key del HTML ──
\newenvironment{clave}{%
  \par\medskip\noindent
  \begin{tabular}{@{}p{0.4em}@{\hspace{8pt}}p{\dimexpr\linewidth-2em\relax}@{}}
  \textcolor{cmorado}{\rule{3pt}{1.1em}} &
}{\end{tabular}\par\medskip}

% ── Caja de advertencia (roja/naranja fuerte) ──
\definecolor{crojo}{HTML}{B3261E}
\newenvironment{advertencia}{%
  \par\medskip\noindent
  \begin{tabular}{@{}p{0.4em}@{\hspace{8pt}}p{\dimexpr\linewidth-2em\relax}@{}}
  \textcolor{crojo}{\rule{3pt}{1.1em}} &
}{\end{tabular}\par\medskip}

\newcommand{\bigsep}{\par\vspace{3pt}}
\newcommand{\dq}[1]{\emph{``#1''}}

\title{\vspace{-1.4cm}\textbf{X5\_alt\_v2 — descriptiva paralela}\\[6pt]
  \large Cómo usar e interpretar \code{X5\_Q1} / \code{X5\_Q2} / \code{X5\_Q3}\\[4pt]
  \normalsize\color{cgris} Análisis estadístico, fórmulas y puente hacia hipótesis de parámetro}
\author{Proyecto Alginvesting}
\date{\today}

\begin{document}
\maketitle
\thispagestyle{empty}

\begin{center}
\begin{minipage}{0.88\linewidth}
\small\color{cgris}
Este documento describe los tres notebooks de \textbf{X5\_alt\_v2} —evolución paralela de
\code{X5\_alt}, que a su vez corre al lado de \code{X5_macro_brain.py} sin reemplazarlo— y
cómo interpretar cada uno de sus resultados. \code{X5\_Q1.ipynb} construye, por activo, una
tabla continua en el tiempo del \emph{estado del activo}; \code{X5\_Q2.ipynb} cuantifica con
evidencia estadística cómo se relaciona cada variable con el precio y el retorno;
\code{X5\_Q3.ipynb} traduce esa evidencia en hipótesis explícitas
\code{parámetro X5 = f(variable)}, sin implementarlas como reglas ni validarlas con
backtesting. Documento rector: \code{docs/plans/x5\_contexto\_conversacion.md}. Estado de
avance citado: 2026-10-12, los tres notebooks completos y ejecutados de punta a punta sobre
BTCUSD.
\end{minipage}
\end{center}

\vspace{8pt}
\tableofcontents
\newpage

% ═══════════════════════════════════════════════════════════════════════════
\section{Introducción}
% ═══════════════════════════════════════════════════════════════════════════

\subsection{Qué es X5\_alt\_v2 y por qué existe}

\code{X5\_macro\_brain.py} (el X5 en producción) funciona, pero acopla demasiados
componentes para iterar con confianza: LightGBM, FT-Transformer, Optuna, ascenso por
gradiente, airbag y un modo demo narrado, todo junto. \textbf{X5\_alt} nació como una
reconstrucción \emph{paralela} —no un reemplazo— para entender primero los datos con
evidencia, antes de decidir qué relaciones vale la pena modelar.

\bigsep
\textbf{X5\_alt\_v2} es, a su vez, una evolución paralela de X5\_alt: no toca
\code{X5\_macro\_brain.py}, \code{X1\_trading.py}, ni \code{config.py}
(\code{TIPO\_EJECUCION}). Tres notebooks en \code{scripts/}, parametrizados por una
\textbf{lista} de activos (\code{valores = ['BTCUSD']}, a diferencia del \code{valor}
singular de la versión anterior):

\begin{itemize}[leftmargin=1.4em,itemsep=2pt]
\item \code{X5\_Q1.ipynb} — tabla de estado del activo (Sección~\ref{sec:q1}).
\item \code{X5\_Q2.ipynb} — core descriptivo: correlación clásica, coseno, no linealidad,
  ceteris paribus, estabilidad, régimen, redundancia, síntesis (Sección~\ref{sec:q2}).
\item \code{X5\_Q3.ipynb} — hipótesis parámetro$\leftrightarrow$variable y tabla final
  (Sección~\ref{sec:q3}).
\end{itemize}

\begin{nota}
\textbf{Documento rector}: \code{docs/plans/x5\_contexto\_conversacion.md} prevalece sobre
cualquier definición de este documento en caso de contradicción. El pre-registro de
hipótesis (anterior a correr cualquier notebook) vive en
\code{docs/plans/X5\_alternativo.md}, sección 10.5 (\code{[PRE-1]}).
\end{nota}

\subsection{Qué cambió respecto a la versión anterior}

La versión anterior (\code{X5\_P1.ipynb}/\code{X5\_P2.ipynb}, hoy congelada como referencia
histórica en \code{scripts/otros\_strings/}) ya construía una tabla continua en el tiempo y
medía correlación clásica. La Tabla~\ref{tab:v1-v2} resume lo que cambió.

\begin{table}[h]
\centering
\small
\begin{tabularx}{\linewidth}{@{}l X X@{}}
\toprule
\textbf{Dimensión} & \textbf{Versión anterior} & \textbf{X5\_alt\_v2} \\
\midrule
Alcance del análisis & podía mezclar drawdown/equity de cuenta & \textbf{siempre por
  activo} — nunca cuenta/posiciones \\
Parámetros \(Y\) a explicar & sin lista cerrada & \code{N}, \code{A}, \code{factor\_B}
  (\(B=\text{factor\_B}\cdot A\)), \code{K}, \code{N\_EXP}, \code{LAMBDA};
  \code{LOTAJES\_M} fijo, \code{PERDIDA\_MAX} pausado \\
Variables \(X\) & indicadores técnicos crudos de X3 & vector de \emph{estado del activo}:
  drawdown, volatilidad, ruido, tendencia/régimen, liquidez, fundamental — más
  transformaciones por ATR en vez de niveles crudos \\
Métrica de asociación & Pearson + Spearman & + \textbf{coseno} \(\cos(P_v,X_v)\)
  (Sección~\ref{ssec:coseno}) \\
Convención de nombres & \code{P1} / \code{P2} & \code{Q1} / \code{Q2} / \code{Q3} \\
Input & \code{valor} (string único) & \code{valores} (lista) \\
Entrega de resultados & archivos sueltos & 1 \code{.zip} por notebook; nunca gráficos
  embebidos — solo código y los datos detrás de cada gráfico, en CSV \\
\bottomrule
\end{tabularx}
\caption{Diferencias entre la versión anterior (congelada) y X5\_alt\_v2.}
\label{tab:v1-v2}
\end{table}

\subsection{Cinco principios que gobiernan los tres notebooks}

\begin{enumerate}[leftmargin=1.4em,itemsep=2pt]
\item \textbf{Parsimonia.} Pocas variables explicativas por parámetro, elegidas por
  evidencia — nunca todas las disponibles.
\item \textbf{Especialización por parámetro.} Cada parámetro depende solo de las variables
  que tengan sentido para él, no de un estado global único.
\item \textbf{Evidencia antes que arquitectura.} No saltar a modelos sofisticados (Markov,
  Aprendizaje por Refuerzo) sin demostrar primero, con datos, que la relación existe.
\item \textbf{Confiabilidad antes que sofisticación.} Entender qué información hay
  disponible y con qué calidad antes de construir nada encima.
\item \textbf{Baseline primero.} Cualquier sofisticación futura debe superar algo simple y
  estático.
\end{enumerate}

\bigsep
A estos se agrega la \textbf{filosofía estadística} que gobierna todo \code{X5\_Q2.ipynb}:
no confundir correlación con causalidad; no confundir relación contemporánea con
predicción (\(\mathrm{Corr}(X_t,P_t)\) no implica anticipar \(P_{t+h}\)); no confundir
repetición de datos con información nueva (los fundamentales llegan por
\emph{forward-fill}); no asumir linealidad ni estabilidad; no sobrecomplicar —favorecer
interpretabilidad sobre maximizar una métrica de \emph{Machine Learning} (ML).

\subsection{Convención de fuerza y dirección}
\label{ssec:convencion-fuerza}

Usada en todo el documento al reportar Pearson (y, por referencia, Spearman/coseno al
lado), heredada tal cual de la versión anterior.

\begin{table}[h]
\centering
\small
\begin{tabular}{@{}l l@{}}
\toprule
\(|\text{Pearson}|\) & Etiqueta \\
\midrule
\(< 0.10\) & nula / despreciable \\
\(0.10 - 0.30\) & débil \\
\(0.30 - 0.50\) & moderada \\
\(0.50 - 0.70\) & fuerte \\
\(\geq 0.70\) & muy fuerte \\
\bottomrule
\end{tabular}
\caption{Convención de fuerza/dirección, usada como referencia incluso al reportar Spearman
y coseno.}
\end{table}

\newpage
% ═══════════════════════════════════════════════════════════════════════════
\section{Notación y glosario de siglas}
\label{sec:notacion}
% ═══════════════════════════════════════════════════════════════════════════

Símbolos usados a lo largo del documento, y las siglas técnicas que aparecen primero en su
forma completa en cada sección pero se abrevian después de la primera mención.

\begin{table}[h]
\centering
\small
\begin{tabularx}{\linewidth}{@{}l X@{}}
\toprule
\textbf{Símbolo} & \textbf{Significado} \\
\midrule
\(P_v\), \(X_v\) & Vector de precios y vector de una variable explicativa, en el tiempo
  (Sección~\ref{ssec:coseno}). \\
\(r_t\) & Retorno logarítmico horario, \(r_t = \ln(\text{Close}_t/\text{Close}_{t-1})\). \\
\(h\) & Horizonte en horas (velas H1) de una ventana retrospectiva o un horizonte de
  retorno pasado. \\
\(\bar{x}\), \(\sigma_x\) & Media y desviación estándar de una variable \(x\). \\
\(X_z\) & Variable estandarizada: \(X_z = (X-\mu_X)/\sigma_X\) (z-score). \\
\(n\), \(k\) & Número de observaciones y número de columnas (regresores + intercepto) de
  una regresión. \\
\(P_0\), \(P_i\) & Precio actual y precio de referencia (soporte o apertura de una
  posición), notación heredada de \code{docs/plans/formulas.tex}. \\
\(A\), \(B\), \(L\) & Umbral de activación en USD, distancia del trailing en USD, y
  exposición en USD por unidad de precio (\(L = \text{lote}\cdot\text{unidades}\)). \\
\bottomrule
\end{tabularx}
\caption{Notación matemática usada en este documento.}
\end{table}

\begin{table}[h]
\centering
\small
\begin{tabularx}{\linewidth}{@{}l X@{}}
\toprule
\textbf{Sigla} & \textbf{Concepto completo} \\
\midrule
ATR & Average True Range (rango verdadero promedio) — medida de volatilidad de Wilder. \\
RSI & Relative Strength Index (índice de fuerza relativa) — oscilador de Wilder, acotado
  0--100. \\
ADX & Average Directional Index (índice direccional promedio) — fuerza de tendencia de
  Wilder, sin dirección. \\
DI\textsuperscript{+}/DI\textsuperscript{-} & Directional Indicator positivo/negativo —
  componentes direccionales que alimentan el ADX. \\
RVOL & Relative Volume (volumen relativo) — \code{Tick\_Volume} contra su propio promedio
  móvil. \\
VIF & Variance Inflation Factor (Factor de Inflación de Varianza) — mide colinealidad
  conjunta de una variable con todas las demás. \\
OLS & Ordinary Least Squares (mínimos cuadrados ordinarios) — método clásico de
  regresión lineal. \\
dcor & Distance Correlation (correlación de distancia), de Székely--Rizzo — captura
  cualquier forma de dependencia, no solo lineal. \\
MI & Mutual Information (información mutua) — dependencia estadística general entre dos
  variables, medida en nats. \\
MAE/MFE & Maximum Adverse/Favorable Excursion (excursión adversa/favorable máxima) — cuánto
  se movió el precio en contra o a favor de una posición mientras estuvo abierta. \\
OE/OA/OC & Orden en Espera / Orden Abierta / Orden Cerrada — progresión de una orden en
  \code{X1\_trading.py} (ver \code{CLAUDE.md}). \\
\bottomrule
\end{tabularx}
\caption{Glosario de siglas técnicas usadas en este documento.}
\end{table}

\newpage
% ═══════════════════════════════════════════════════════════════════════════
\section{X5\_Q1 — tabla de estado del activo}
\label{sec:q1}
% ═══════════════════════════════════════════════════════════════════════════

\subsection{Arquitectura y pipeline}

Por cada activo de \code{valores}, el notebook construye una tabla H1 continua (una fila
por vela horaria): precio + X3 (técnicos) + X2 (fundamental) + seis familias de variables de
\emph{estado del activo}, en este orden —cada sección encadena sobre la tabla de la
anterior—:

\begin{center}
\small
\begin{tabular}{@{}cccccccc@{}}
base & \(\to\) & técnicos & \(\to\) & drawdown & \(\to\) & volatilidad & \(\to\) \\
liquidez & \(\to\) & fundamental & \(\to\) & ruido & \(\to\) & tendencia/régimen & \(\to\) \\
\multicolumn{8}{c}{catálogo de documentación \(\to\) guardar CSV \(\to\) exportar \code{.zip}} \\
\end{tabular}
\end{center}

Guarda \code{\{valor\}\_tabla\_estado\_v2.csv} y \code{\{valor\}\_catalogo\_variables.csv}
en \code{resources/x5\_alt/} — consumidos por Q2/Q3, nunca recalculados ahí.
\textbf{Reutiliza en vez de reimplementar}: \code{\_calcular\_todos\_indicadores} y
\code{\_wilder\_ema} de \code{X3\_technical\_features.py} — ningún indicador (RSI, ATR,
ADX) se recalcula desde cero donde X3 ya lo resuelve.

\subsection{Cómo correr el notebook}

\begin{enumerate}[leftmargin=1.4em,itemsep=2pt]
\item Abrir \code{scripts/X5\_Q1.ipynb} en Windows (necesita \code{Data/},
  \code{resources/x2/} y \code{resources/x3/} actualizados — no disponibles así en Mac).
\item Editar, si corresponde, la primera celda de código: \code{valores = ['BTCUSD']} —
  es una \textbf{lista}, se puede extender a varios activos sin tocar el resto del
  notebook.
\item Ejecutar todas las celdas de arriba hacia abajo (\code{Run All} o
  \code{jupyter nbconvert --execute}, kernel \code{revenAI}). Cada sección depende de la
  anterior — no hay celdas que se puedan saltar.
\item Verificar al final que se generaron, por activo,
  \code{\{valor\}\_tabla\_estado\_v2.csv} y \code{\{valor\}\_catalogo\_variables.csv} en
  \code{resources/x5\_alt/}, y que la última celda empaquetó el \code{.zip} con
  \code{exportar\_zip}.
\item Bajar el \code{.zip} de Windows y pasarlo a Claude en Mac — es la vía normal de
  entrega, ya que \code{Data/}/X2/X3 no están disponibles así en Mac.
\end{enumerate}

\subsection{Técnicos filtrados y transformados}

Los niveles crudos de medias móviles (\code{x3\_sma\_\{20,50,200\}},
\code{x3\_ema\_\{12,26\}}) se descartan: su correlación con \code{Close} es \(\approx 1\)
por construcción, no aportan señal nueva como variable de estado. Se reemplazan por dos
transformaciones normalizadas por ATR —independientes de la escala de precio y del régimen
de volatilidad—:

\begin{align}
\text{distancia\_sma}_w &= \frac{\text{Close} - \text{SMA}_w}{\text{ATR}_{14}} \\
\text{pendiente\_sma}_w &= \frac{\text{SMA}_w.\text{diff}(w) / w}{\text{ATR}_{14}}
\end{align}

con \(w \in \{20,50,200\}\) para SMA y \(w \in \{12,26\}\) para EMA (misma fórmula,
sustituyendo \(\text{SMA}_w\) por \(\text{EMA}_w\)).

\subsection{Variables de drawdown del activo}

Drawdown \textbf{del activo} —nunca de cuenta o posiciones, eso es responsabilidad de
X1/X4—, en dos vistas complementarias:

\begin{align}
\text{drawdown\_historico} &= \frac{\text{Close} - \mathrm{cummax}(\text{Close})}
  {\mathrm{cummax}(\text{Close})} \\
\text{duracion\_drawdown} &= \text{velas desde el último máximo histórico} \\
\text{recuperacion\_desde\_minimo} &= \frac{\text{Close} - \min(\text{episodio})}
  {\min(\text{episodio})} \\
\text{velocidad\_caida} &= \frac{\text{drawdown\_historico}}{\text{duracion\_drawdown}} \\
\text{drawdown}_{d} &= \frac{\text{Close} - \max_{h}(\text{Close})}{\max_{h}(\text{Close})},
  \quad d \in \{7,30,90,365\}\text{ días}
\end{align}

donde \(h\) es la ventana rodante en horas (\(168, 720, 2160, 8760\)). A diferencia del
histórico, las ventanas de \(d\) días \dq{resetean} la referencia si el precio alcanzó un
máximo reciente, aunque no sea el máximo histórico absoluto.

\subsubsection*{Ejemplo numérico (ilustrativo, no datos reales)}

Diez velas de juguete, para ver cómo se mueven \code{drawdown\_historico},
\code{duracion\_drawdown}, \code{velocidad\_caida} y \code{recuperacion\_desde\_minimo}
juntas.

\begin{table}[h]
\centering
\small
\begin{tabular}{@{}r r r r r r r@{}}
\toprule
\(t\) & Close & cummax & drawdown\_hist. & duracion & velocidad\_caida & recuperación \\
\midrule
1  & 100 & 100 & 0{,}000  & 0 & — & 0{,}000 \\
2  & 102 & 102 & 0{,}000  & 0 & — & 0{,}000 \\
3  & 101 & 102 & \(-0{,}010\) & 1 & \(-0{,}0098\) & 0{,}000 \\
4  & 105 & 105 & 0{,}000  & 0 & — & 0{,}000 \\
5  & 103 & 105 & \(-0{,}019\) & 1 & \(-0{,}0190\) & 0{,}000 \\
6  & 98  & 105 & \(-0{,}067\) & 2 & \(-0{,}0333\) & 0{,}000 \\
7  & 99  & 105 & \(-0{,}057\) & 3 & \(-0{,}0190\) & 0{,}010 \\
8  & 104 & 105 & \(-0{,}010\) & 4 & \(-0{,}0024\) & 0{,}061 \\
9  & 106 & 106 & 0{,}000  & 0 & — & 0{,}000 \\
10 & 103 & 106 & \(-0{,}028\) & 1 & \(-0{,}0283\) & 0{,}000 \\
\bottomrule
\end{tabular}
\caption{Serie de juguete. \dq{velocidad\_caida} queda sin definir (\code{NaN}) cuando
\code{duracion\_drawdown}\(=0\) —justo en la vela donde se marca un nuevo máximo— por la
división por cero evitada con \code{replace(0, NaN)}.}
\end{table}

Lectura: en \(t{=}6\) el precio cae a 98, lejos del máximo de 105 marcado en \(t{=}4\)
—drawdown de \(-6{,}7\%\) en solo 2 velas, la caída más rápida de la serie
(\(\text{velocidad\_caida}=-0{,}0333\))—. Entre \(t{=}7\) y \(t{=}8\) el precio se recupera
desde el mínimo del episodio (98) sin alcanzar aún un nuevo máximo histórico, por eso
\code{recuperacion\_desde\_minimo} crece (\(0{,}010 \to 0{,}061\)) mientras
\code{drawdown\_historico} sigue negativo. Recién en \(t{=}9\) el precio supera 105 y
\code{duracion\_drawdown} vuelve a 0 — comienza un episodio nuevo.

\subsection{Variables de volatilidad}

El ATR absoluto/relativo y la volatilidad realizada corta/larga ya los calcula X3
(\code{x3\_atr\_14}, \code{x3\_atr\_pct}, \code{x3\_vol\_24h}, \code{x3\_vol\_7d}) — se
reutilizan tal cual. Se agregan tres variables nuevas, sobre retornos logarítmicos
horarios \(r_t = \ln(\text{Close}_t / \text{Close}_{t-1})\):

\begin{align}
\text{cambio\_volatilidad} &= \frac{x3\_vol\_24h}{x3\_vol\_7d} \\
\text{vol\_alcista}_{30d} &= \mathrm{std}_{30d}(r_t \mid r_t > 0) \cdot \sqrt{8760} \\
\text{vol\_bajista}_{30d} &= \mathrm{std}_{30d}(r_t \mid r_t < 0) \cdot \sqrt{8760} \\
\text{frecuencia\_extremos}_{30d} &= \mathrm{media}_{30d}\big(|r_t| > 2\cdot
  \mathrm{std}_{\mathrm{expandida}}(r_t)\big)
\end{align}

\begin{nota}
El umbral de \code{frecuencia\_extremos\_30d} usa \code{expanding()} (no \code{rolling()})
para no mirar información futura al definirlo — el umbral en \(t\) solo usa información
\(\leq t\).
\end{nota}

\subsection{Variables de liquidez}

\code{Spread} y \code{Tick\_Volume} vienen directo del CSV de MT5 (\code{Real\_Volume}
siempre en 0 — el bróker no lo reporta).

\begin{align}
\text{spread\_relativo} &= \frac{\text{Spread}}{\mathrm{media}_{30d}(\text{Spread})} \\
\text{tick\_volume\_z} &= \frac{\text{Tick\_Volume} - \mathrm{media}_{30d}}
  {\mathrm{std}_{30d}} \\
\text{volumen\_relativo} &= \frac{\text{Tick\_Volume}}{\mathrm{media}_{30d}
  (\text{Tick\_Volume})}
\end{align}

\begin{nota}
\code{Spread} viene en puntos del bróker MT5, no en USD — convertirlo exige
\code{symbol\_info(valor).point}, no disponible en Mac ni persistido en los CSV. Por eso
\code{spread\_relativo} normaliza contra su propio promedio móvil en vez de contra
\code{Close}.
\end{nota}

\subsection{Variables fundamentales de estado}

Las columnas \code{x2\_score}/\code{x2\_score\_cross}/\code{x2\_score\_tendencia} llegan
por \code{merge\_asof(direction='backward')} —forward-fill del último valor conocido—.
Las variables de esta sección existen para distinguir una actualización real de la mera
repetición del forward-fill:

\begin{align}
\text{fundamental\_cambio} &= x2\_score.\mathrm{diff}() \\
\text{fundamental\_edad} &= \text{velas desde la última actualización real} \\
\text{fundamental\_contradiccion} &= (x2\_score\_cross - 0.5)\cdot
  (x2\_score\_tendencia - 0.5)
\end{align}

En \code{fundamental\_contradiccion}: negativo = contradicción (un componente sobre
neutro y el otro bajo), positivo = consistencia, \(0\) = al menos uno es neutro.

\begin{advertencia}
\textbf{\(x2\_score\) queda 100\% \texttt{NaN} para BTCUSD con los datos actuales de Mac.}
\code{Data/BTCUSD.csv} termina el 2026-06-02; los únicos registros de
\code{x2\_history.json} para BTCUSD son posteriores (2026-06-12+) → el \emph{merge asof
backward} no tiene nada hacia atrás que propagar. \textbf{No es un bug}: es la escasez de
X2 ya señalada en \code{X5\_alternativo.md}, agravada por el desfase de fechas en esta
copia. Se valida al refrescar \code{Data/}/\code{x2\_history.json} desde Windows con fechas
que se solapen.
\end{advertencia}

\subsection{Variables de ruido}

Sobre retornos logarítmicos horarios y ventana de 30 días (720h), reutilizando
\code{x3\_atr\_pct} en vez de recalcular ATR:

\begin{align}
\text{ruido\_cambios\_signo}_{30d} &= \mathrm{media}_{30d}\big(\mathrm{sign}(r_t) \neq
  \mathrm{sign}(r_{t-1})\big) \\
\text{ruido\_ratio\_movimiento}_{30d} &= \frac{\sum_{30d} |r_t|}
  {\left|\sum_{30d} r_t\right|} \\
\text{ruido\_mecha}_{30d} &= \mathrm{media}_{30d}\left(
  \frac{\text{High}-\text{Low}-|\text{Close}-\text{Open}|}{\text{High}-\text{Low}}\right) \\
\text{ruido\_atr\_sin\_direccion}_{30d} &= \frac{\mathrm{media}_{30d}(x3\_atr\_pct)}
  {\left|\sum_{30d} r_t\right|}
\end{align}

\begin{nota}
\code{ruido\_ratio\_movimiento\_30d} es el \textbf{inverso del Efficiency Ratio de
Kaufman}: cercano a 1 = recorrido \(\approx\) desplazamiento neto (movimiento eficiente);
alto = mucho recorrido para poco desplazamiento neto (ruido/\emph{whipsaw}).
\code{ruido\_mecha\_30d} agrega en ventana el mismo concepto de \dq{fuerza del rechazo} que
usa el algoritmo de soportes de X0 (\(f = 1-|\text{Close}-\text{Open}|/(\text{High}-
\text{Low})\)) por vela.
\end{nota}

\subsection{Variables de tendencia y régimen}

\begin{align}
\mathrm{ADX}_{14} &= \mathrm{WilderEMA}(\mathrm{DX}, 14), \qquad
  \mathrm{DX} = 100 \cdot \frac{|{+}\mathrm{DI} - {-}\mathrm{DI}|}{{+}\mathrm{DI}+{-}
  \mathrm{DI}}
\end{align}

reutilizando \code{\_wilder\_ema} —el mismo suavizado que \code{\_calcular\_atr}— en vez de
un esquema de suavizado distinto.

\begin{lstlisting}
regimen_tendencia:
    'sideways'  si adx_14 < 20                 # mercado sin tendencia clara (Wilder)
    'bull'      si adx_14 >= 20 y pendiente > 0  # pendiente = x3_trend_slope_50
    'bear'      si adx_14 >= 20 y pendiente < 0

ruptura_rango:
    +1  si Close > max(20 velas previas)       # estilo canal de Donchian
    -1  si Close < min(20 velas previas)
     0  en otro caso

duracion_en_rango = velas desde la última ruptura de rango
\end{lstlisting}

\subsubsection*{Ejemplo numérico simplificado (ilustrativo, no datos reales)}

Cuatro velas de juguete, para ver de dónde salen \(+\mathrm{DM}\), \(-\mathrm{DM}\),
\(\mathrm{TR}\), \(\pm\mathrm{DI}\) y \(\mathrm{DX}\) \textbf{antes} del suavizado de
Wilder (que en producción corre sobre 14 períodos, no 3 — acá se usa un promedio simple
solo para ilustrar la mecánica, no el ADX real).

\begin{table}[h]
\centering
\small
\begin{tabular}{@{}r r r r r r r@{}}
\toprule
\(t\) & High & Low & Close\(_{prev}\) & \(+\mathrm{DM}\) & \(-\mathrm{DM}\) & TR \\
\midrule
2 & 104 & 101 & 100 & 3 & 0 & 4 \\
3 & 103 & 99  & 103 & 0 & 2 & 4 \\
4 & 107 & 102 & 100 & 4 & 0 & 7 \\
\bottomrule
\end{tabular}
\caption{\(+\mathrm{DM}\)/\(-\mathrm{DM}\)/TR de la serie de juguete (High inicial 101,
Low inicial 99 en \(t{=}1\), sin fila propia por no tener vela previa).}
\end{table}

Promediando simple las 3 filas (en vez del suavizado de Wilder de 14 períodos):

\begin{align}
\overline{\mathrm{TR}} = \frac{4+4+7}{3} = 5, \qquad
\overline{{+}\mathrm{DM}} = \frac{3+0+4}{3} = 2{,}33, \qquad
\overline{{-}\mathrm{DM}} = \frac{0+2+0}{3} = 0{,}67
\end{align}

\begin{align}
{+}\mathrm{DI} = 100 \cdot \frac{2{,}33}{5} = 46{,}7, \qquad
{-}\mathrm{DI} = 100 \cdot \frac{0{,}67}{5} = 13{,}3
\end{align}

\begin{align}
\mathrm{DX} = 100 \cdot \frac{|46{,}7 - 13{,}3|}{46{,}7 + 13{,}3} = 100 \cdot
  \frac{33{,}3}{60{,}0} \approx 55{,}6
\end{align}

Si este nivel se sostuviera tras el suavizado de 14 períodos, \(\mathrm{ADX}_{14}
\approx 55{,}6 > 20\) \(\Rightarrow\) \textbf{no} es \code{'sideways'}; como
\(+\mathrm{DM}\) domina sobre \(-\mathrm{DM}\) (la pendiente de \code{x3\_trend\_slope\_50}
sería positiva en un contexto así), \code{regimen\_tendencia} resultaría \code{'bull'}.

\subsection{Catálogo de variables y limitaciones}

La última sección del notebook genera, por columna de la tabla final (incluidas las
heredadas de precio/X3/X2), una fila con fuente, significado, unidad, fórmula,
\code{forward\_fill}, granularidad esperada/observada, cobertura y limitaciones —
referencia rápida para \code{X5\_Q2}/\code{X5\_Q3} sin releer cada celda de este notebook.
Se guarda en \code{resources/x5\_alt/\{valor\}\_catalogo\_variables.csv}.

\bigsep
\textbf{Limitaciones adicionales}, más allá del \(x2\_score\) ya señalado: NaN de
warm-up en toda variable con ventana (ej. \code{x3\_vol\_7d} en las primeras 167 velas;
\code{drawdown\_\{7,30,90,365\}d} usan \code{min\_periods=1} desde la primera vela, a
diferencia de \code{x3\_drawdown\_20/50} que exigen la ventana completa).

\newpage
% ═══════════════════════════════════════════════════════════════════════════
\section{X5\_Q2 — core descriptivo}
\label{sec:q2}
% ═══════════════════════════════════════════════════════════════════════════

Segundo notebook — consume \code{\{valor\}\_tabla\_estado\_v2.csv} de Q1, nunca la
recalcula. Cubre diez ítems analíticos (\code{Q2-2} a \code{Q2-14}), todos sobre el mismo
conjunto de variables explicativas y con el mismo filtro mínimo: 30 observaciones no nulas
y ambas series con dispersión no nula; si no se cumple, se omite con aviso.

\subsection{Cómo correr el notebook}

\begin{enumerate}[leftmargin=1.4em,itemsep=2pt]
\item Requiere que \code{X5\_Q1.ipynb} ya haya generado
  \code{\{valor\}\_tabla\_estado\_v2.csv} para cada activo de \code{valores} — si falta,
  la primera celda de carga levanta \code{FileNotFoundError} con un mensaje explícito.
\item Ejecutar de arriba hacia abajo. Las catorce subsecciones (\code{Q2-2} a
  \code{Q2-14}) son mayormente independientes entre sí, pero todas dependen de la carga
  inicial (\code{Q2-1}) y comparten las mismas variables/targets definidos ahí
  (\code{\_variables\_analisis\_q2\_3}, \code{TARGETS\_Q2\_3}) — no se recalculan por
  separado en cada sección.
\item Es el notebook más largo de los tres (89 celdas) — correrlo completo en Windows
  toma varios minutos, principalmente por la correlación de distancia (\code{Q2-14},
  cuadrática en el tamaño de la submuestra) y los 14 ítems con sus gráficos.
\item Verificar al final que se generó
  \code{\{valor\}\_q2\_5\_export\_metricas.json} —es el único archivo que
  \code{X5\_Q3.ipynb} necesita de este notebook— y que el \code{.zip} final incluye los
  \textasciitilde19 CSV intermedios (uno por ítem, por activo).
\end{enumerate}

\subsection{Decisión de diseño: dos paneles, no uno}
\label{ssec:dos-paneles}

Cada variable se evalúa contra \textbf{dos targets}, nunca solo contra \code{Close}:

\begin{itemize}[leftmargin=1.4em,itemsep=2pt]
\item \textbf{Panel A — nivel (\code{Close}), referencial.} Fiel al planteamiento
  original (\dq{variable$\leftrightarrow$precio}), se muestra pero se marca expresamente como
  \textbf{no concluyente}: \code{Close} no es estacionario (BTCUSD recorrió
  \(\sim\)20.000\(\to\)100.000 USD entre 2021 y 2026), así que casi cualquier variable con
  tendencia propia \dq{correlaciona} con el nivel por coincidencia temporal.
\item \textbf{Panel B — retorno contemporáneo, principal.}
  \(\text{retorno\_log\_1h} = \ln(\text{Close}_t / \text{Close}_{t-1})\). Estacionario —
  alimenta el ranking de esta sección y todos los ítems posteriores.
\end{itemize}

\code{regimen\_tendencia} (bull/bear/sideways, sin orden natural) se trata aparte:
comparación directa de medias por categoría en vez de forzar un \code{factorize}
arbitrario en Pearson/Spearman.

\subsection{Pearson, Spearman y decil (Q2-3)}

Por variable y panel: Pearson (relación lineal), Spearman (relación monótona), scatter +
regresión lineal simple (descriptiva, no predictiva), y el target medio por decil (10
cuantiles) de la variable.

\begin{align}
r_{\mathrm{Pearson}} &= \frac{\sum_i (x_i - \bar{x})(y_i - \bar{y})}
  {\sqrt{\sum_i (x_i-\bar{x})^2} \cdot \sqrt{\sum_i (y_i-\bar{y})^2}} \\
\rho_{\mathrm{Spearman}} &= r_{\mathrm{Pearson}}\big(\mathrm{rango}(x),\, \mathrm{rango}(y)
  \big) \;=\; 1 - \frac{6\sum_i d_i^2}{n(n^2-1)} \quad \text{(sin empates)}
\end{align}

donde \(d_i = \mathrm{rango}(x_i) - \mathrm{rango}(y_i)\). Spearman es, literalmente,
Pearson aplicado a los rangos en vez de a los valores crudos — por eso detecta relaciones
monótonas no lineales que Pearson puede subestimar.

\subsubsection*{Ejemplo numérico (ilustrativo)}

Con \(x=(1,2,3,4,5)\) y \(y=(2,1,4,3,5)\) —monótona pero no perfectamente—:
\(\bar{x}=\bar{y}=3\); \(\sum(x_i-\bar{x})(y_i-\bar{y}) = 8\);
\(\sum(x_i-\bar{x})^2 = \sum(y_i-\bar{y})^2 = 10\) \(\Rightarrow\)

\begin{align}
r_{\mathrm{Pearson}} = \frac{8}{\sqrt{10\cdot 10}} = 0{,}80
\end{align}

Como \(x\) ya son sus propios rangos y \(y\) es una permutación de \(1..5\), los rangos de
\(y\) son \(y\) mismo: \(d = (1{-}2,\,2{-}1,\,3{-}4,\,4{-}3,\,5{-}5) = (-1,1,-1,1,0)\),
\(\sum d_i^2 = 4\):

\begin{align}
\rho_{\mathrm{Spearman}} = 1 - \frac{6\cdot 4}{5\cdot 24} = 1 - 0{,}20 = 0{,}80
\end{align}

En este caso particular ambas coinciden (\(0{,}80\), \dq{fuerte} según la convención de la
Sección~\ref{ssec:convencion-fuerza}) porque los datos de juguete ya eran casi rangos puros
— con datos reales ruidosos, Pearson y Spearman típicamente divergen más, y esa divergencia
es justamente lo que revela si la relación es lineal o solo monótona.

\begin{lstlisting}
def analizar_variable_vs_target(df, col, target_col, n_bins=10):
    sub = df[[col, target_col]].dropna()
    if len(sub) < 30: return None
    pear_r, pear_p   = pearsonr(x, y)
    spear_r, spear_p = spearmanr(x, y)
    bins = pd.qcut(sub[col], n_bins, duplicates='drop')
    por_decil = sub.groupby(bins)[target_col].mean()
\end{lstlisting}

\subsubsection*{Cómo leer la tabla de deciles}

\code{pd.qcut} divide los valores de la variable en 10 grupos de igual tamaño
—\emph{cuantiles}, no rangos fijos—: el decil 1 son el 10\% de las observaciones con el
valor \textbf{más bajo} de la variable, el decil 10 el 10\% con el valor \textbf{más
alto}, y cada decil tiene el mismo número de observaciones por construcción (a diferencia
de dividir el rango de la variable en 10 tramos iguales, que dejaría deciles con
cantidades de datos muy distintas si la variable no es uniforme). La columna que importa
es el \textbf{target medio dentro de cada decil}: si esos 10 promedios crecen de forma
pareja del decil 1 al 10, la relación es aproximadamente lineal y monótona; si crecen
solo en los extremos y quedan planos en el medio, es la señal de un posible umbral (ver
\dq{no linealidades}, más adelante en esta sección); si no siguen ningún patrón claro, la
variable probablemente no tiene relación real con el target más allá de lo que ya
capturó Pearson.

\textbf{Limitación central}: \(\sim\)40.000 observaciones horarias fuertemente
autocorrelacionadas (no independientes) \(\Rightarrow\) los \(p\)-values salen
extremadamente pequeños incluso para relaciones triviales — se muestran solo de
referencia, no deben leerse literalmente.

\subsection{Coseno vector-precio vs. vector-variable (Q2-4)}
\label{ssec:coseno}

Complementa Pearson/Spearman con una medida de \textbf{alineación direccional} entre el
vector de precios \(P_v\) y el de cada variable \(X_v\) en el tiempo:

\begin{align}
\cos(P_v, X_v) &= \frac{P_v \cdot X_v}{\|P_v\|\,\|X_v\|}
\end{align}

A diferencia de Pearson, \textbf{no centra nada} — mide el ángulo entre los vectores
crudos. Pearson \emph{es} el coseno de las series ya centradas:

\begin{align}
\mathrm{corr}(x,y) = \cos(x - \bar{x},\; y - \bar{y})
\end{align}

Esa diferencia produce comportamientos distintos según el panel, verificados con BTCUSD
real (Tabla~\ref{tab:coseno-ejemplo}).

\begin{table}[h]
\centering
\small
\begin{tabularx}{\linewidth}{@{}X c c c X@{}}
\toprule
\textbf{Variable} & \textbf{Panel} & \textbf{Pearson} & \textbf{Coseno} & \textbf{Lectura} \\
\midrule
\code{drawdown\_historico} & A (\code{Close}) & \(+0.87\) & \(-0.51\) & el coseno sin
  centrar \textbf{invierte el signo} frente a Pearson \\
\code{ruido\_cambios\_signo\_30d} & A (\code{Close}) & \(-0.54\) & \(+0.87\) & mismo
  fenómeno, en sentido opuesto \\
\code{x3\_rsi\_14} (centrado \(\sim\)50) & B (retorno) & \(0.338\) & \(0.080\) & coseno
  fuertemente atenuado por el baseline alto \\
\code{x3\_macd\_hist} (centrado \(\sim\)0) & B (retorno) & \(0.167\) & \(0.167\) & sin
  baseline que distorsione, coinciden casi exacto \\
\bottomrule
\end{tabularx}
\caption{Coseno vs. Pearson, ejemplos reales de BTCUSD. El coseno no es \dq{Pearson más
permisivo} sobre el nivel — es una cantidad distinta.}
\label{tab:coseno-ejemplo}
\end{table}

\subsection{No linealidades, umbrales y saturación (Q2-7)}

\textbf{Pregunta}: ¿la relación variable$\leftrightarrow$target es lineal en todo su rango, o cambia de
forma en algún tramo —plana en valores normales pero distinta sobre el percentil 90, o con
un efecto que se satura en los extremos—?

\begin{align}
\Delta R^2 &= R^2_{\mathrm{cuad}} - R^2_{\mathrm{lin}}
  \qquad (\texttt{polyfit} \text{ grado 1 vs. 2})
\end{align}

\begin{table}[h]
\centering
\small
\begin{tabularx}{\linewidth}{@{}l X l@{}}
\toprule
\textbf{Métrica} & \textbf{Regla} & \textbf{Umbral} \\
\midrule
Curvatura & \(\Delta R^2 \geq 0.01\) \(\Rightarrow\) curvatura relevante; forma U/U
  invertida según signo del coef. cuadrático & \(0.01\) \\
Umbral de cola & \(|\text{media}_{\text{cola}} - \text{media}_{\text{cuerpo}}| > k \cdot
  \text{std}_{\text{cuerpo}}\), colas en \(p_{10}\)/\(p_{90}\), mín. 10 obs. & \(k=0.5\) \\
Saturación & pendiente local de la cola mantiene signo global pero
  \(|\text{pendiente}_{\text{cola}}/\text{pendiente}_{\text{global}}| < r\) & \(r=0.3\) \\
Asimetría & \(\mathrm{skew}(\text{variable})\) + Pearson(mitad baja) vs.
  Pearson(mitad alta) de la mediana & cambio de signo \(\Rightarrow\) no simétrica \\
\bottomrule
\end{tabularx}
\caption{Umbrales exactos de Q2-7.}
\end{table}

\begin{clave}
Una variable con \(\Delta R^2\) alto, cola marcada como umbral, o asimetría fuerte entre
mitades es candidata a jugar un rol de \dq{régimen}/condición —activa o modula un
parámetro solo bajo cierto percentil— en vez de una relación lineal continua. Alimenta las
hipótesis de \code{K}/\code{N\_EXP} (Sección~\ref{sec:q3}).
\end{clave}

\subsection{Variable vs. comportamiento histórico reciente del precio (Q2-8)}

A diferencia de Q2-3 (contemporáneo), aquí el target es lo que el precio ya venía
acumulando \textbf{antes} de \(t\):

\begin{align}
\text{retorno\_pasado}_h = \ln\!\left(\frac{\text{Close}_t}{\text{Close}_{t-h}}\right),
  \qquad h \in \{6h,\, 24h,\, 72h,\, 168h\}
\end{align}

Reutiliza exactamente \code{analizar\_variable\_vs\_target} de Q2-3 —solo cambia el
target—. Sigue siendo descriptivo/histórico: una variable fuertemente asociada al retorno
pasado de 7 días puede estar simplemente \emph{midiendo} ese mismo retorno (variables
construidas desde el propio precio, ej. \code{x3\_roc\_20}) — la redundancia (Q2-12,
Sección~\ref{ssec:redundancia}) distingue eso de una relación informativa nueva.

\subsection{Ceteris paribus: regresión multivariable y VIF (Q2-6)}

\textbf{Pregunta}: si comparo dos instantes con el \textbf{mismo} valor en todas las demás
variables y que solo difieran en una, ¿cambia el precio/retorno? Varias variables de Q1
comparten información —un Pearson individual alto puede ser la misma señal contada dos
veces—.

\begin{align}
X_z &= \frac{X - \mu_X}{\sigma_X}, \qquad y_z = \frac{y - \bar{y}}{\sigma_y} \\
\hat{\beta} &= (X_z'X_z)^{-1}X_z'y_z \qquad \text{(mínimos cuadrados, con intercepto)} \\
\hat{\sigma}^2 &= \frac{\hat{e}'\hat{e}}{n-k}, \qquad \hat{e} = y_z - X_z\hat{\beta} \\
\mathrm{SE}(\hat\beta) &= \sqrt{\mathrm{diag}\!\left(\hat\sigma^2 (X_z'X_z)^{-1}\right)} \\
t &= \frac{\hat\beta}{\mathrm{SE}(\hat\beta)}, \qquad p = 2\cdot P(T > |t|,\; \mathrm{dof}=n-k)
\end{align}

El coeficiente \(\hat\beta_j\) es el cambio esperado en el target (en desvíos estándar)
cuando la variable \(j\) sube 1 desvío estándar, con las demás fijas. Calculado a mano con
\code{numpy}/\code{scipy} (mínimos cuadrados clásicos de OLS), mismo enfoque de
\code{Otros/scripts/X5\_analisis\_exploratorio.ipynb}, adaptado —no importado literal, para
no acoplar X5\_alt\_v2 a la implementación de X5 original.

\bigsep
\textbf{Multicolinealidad} — Factor de Inflación de Varianza (VIF):

\begin{align}
\mathrm{VIF}_j = \mathrm{diag}\!\left(\mathrm{pinv}\big(\mathrm{corr}(X_z)\big)\right)_j
\end{align}

equivalente a \(1/(1-R^2)\) de regresar la variable \(j\) contra todas las demás, sin
correr \(N\) regresiones. \(\mathrm{VIF}>10\) señala colinealidad alta —el coeficiente
individual es poco confiable, sin que eso implique eliminar la variable—.

\begin{nota}
\textbf{Sensibilidad a especificaciones}: se repite la regresión excluyendo las variables
con \(\mathrm{VIF}>10\) (especificación reducida); si el coeficiente de una variable que
sobrevive cambia de \textbf{signo} entre la especificación completa y la reducida, el
resultado depende de qué colineales se incluyen — no es robusto.
\end{nota}

\subsubsection*{Ejemplo numérico de VIF (ilustrativo)}

Tres variables de juguete con esta matriz de correlación —las dos primeras pensadas como
dos medidas de volatilidad muy parecidas (ej. \code{x3\_vol\_24h} y \code{x3\_atr\_pct},
Pearson real \(0{,}91\) en BTCUSD), la tercera independiente—:

\begin{align}
\mathrm{corr}(X_z) = \begin{pmatrix} 1 & 0{,}95 & 0{,}30 \\ 0{,}95 & 1 & 0{,}20 \\
  0{,}30 & 0{,}20 & 1 \end{pmatrix}
\end{align}

Invirtiendo esta matriz \(3\times 3\) y tomando su diagonal (\(\mathrm{VIF}_j =
\mathrm{diag}(\mathrm{pinv}(\mathrm{corr}))_j\)):

\begin{center}
\small
\begin{tabular}{@{}l r@{}}
\toprule
Variable & VIF \\
\midrule
1 (ej. \code{x3\_vol\_24h}) & \textbf{11{,}78} \\
2 (ej. \code{x3\_atr\_pct}) & \textbf{11{,}17} \\
3 (independiente) & 1{,}20 \\
\bottomrule
\end{tabular}
\end{center}

Las variables 1 y 2 cruzan el umbral \(\mathrm{VIF}>10\): su alta correlación mutua
(\(0{,}95\)) hace que sus coeficientes individuales en la regresión completa sean poco
confiables —no se sabe si el efecto es de la variable 1, de la 2, o de ambas a la vez—,
mientras que la variable 3, casi sin colinealidad, mantiene un VIF cercano a 1. Es
exactamente el criterio que usa \code{Q3-7} (Sección~\ref{sec:q3}) para descartar
\code{x3\_vol\_24h} como candidata de \code{A} en favor de \code{x3\_atr\_pct}: agregar
ambas no suma información nueva, solo infla la incertidumbre de los coeficientes.

\subsection{Estabilidad temporal: por año, rolling y coseno por ventanas (Q2-9 / Q2-10)}

Un Pearson global fuerte puede estar promediando sub-períodos con signos opuestos.

\begin{itemize}[leftmargin=1.4em,itemsep=2pt]
\item \textbf{Q2-9 — por año y rolling de 90 días}: Pearson/Spearman/coseno por año
  calendario; además un Pearson \emph{rolling} de 90 días (2.160h), resumido en su
  desviación estándar en el tiempo.
\item \textbf{Q2-10 — coseno rolling por ventanas (30/90/180 días)}: no existe
  \code{.rolling().corr()} para coseno, así que se calcula vectorizado:
\end{itemize}

\begin{align}
\cos_t = \frac{\sum_{\text{ventana}} x\cdot y}{\sqrt{\sum_{\text{ventana}} x^2} \cdot
  \sqrt{\sum_{\text{ventana}} y^2}}
\end{align}

\begin{table}[h]
\centering
\small
\begin{tabular}{@{}l l@{}}
\toprule
\textbf{Clasificación} & \textbf{Condición} \\
\midrule
cambio estructural & el signo cambia entre años/ventanas (o el rolling cruza cero) \\
inestable & mismo signo, dispersión (std) \(> 0.15\) \\
estable & mismo signo y dispersión \(\leq 0.15\) \\
\bottomrule
\end{tabular}
\caption{Clasificación de estabilidad, misma regla en Q2-9 y Q2-10.}
\end{table}

Q2-10 agrega una pregunta propia: \code{consistente\_entre\_ventanas} — ¿la clasificación
es la misma en las 3 ventanas? Si no, la lectura de estabilidad depende de una elección
arbitraria de ventana.

\subsection{Redundancia entre variables explicativas (Q2-12)}
\label{ssec:redundancia}

A diferencia de todo lo anterior (variable$\leftrightarrow$target), mide relaciones \textbf{entre las
propias variables explicativas}: matriz de Pearson entre todas las variables; un par se
marca \textbf{redundante} si \(|\text{Pearson}| > 0.85\) —más estricto que el
\(\mathrm{VIF}>10\) de Q2-6, porque acá se mide similitud bivariada directa entre dos
variables, no el efecto conjunto de todas las demás sobre una—.

\bigsep
La familia técnico/fundamental se deriva de la columna \code{fuente} del catálogo de Q1
(mapeo documentado, auditable). Un par \code{redundante} \textbf{no implica eliminar} una
de las dos variables —solo que, si se usan ambas como explicativas de un mismo parámetro,
aportan señal muy parecida—. Usado en la Sección~\ref{sec:q3} para no elegir dos
candidatas redundantes para el mismo parámetro: ej. \code{x3\_vol\_24h} se descarta frente
a \code{x3\_atr\_pct} (Pearson \(0.91\) entre ambas).

\subsection{Análisis por regímenes (Q2-13)}

Tres dimensiones, cada una con categorías explícitas:

\begin{table}[h]
\centering
\small
\begin{tabularx}{\linewidth}{@{}l l X@{}}
\toprule
\textbf{Dimensión} & \textbf{Categorías} & \textbf{Definición} \\
\midrule
Tendencia & bull / bear / sideways & \code{regimen\_tendencia}, ya calculada en Q1 \\
Volatilidad & alta / baja & split por \textbf{mediana} de \code{x3\_vol\_24h} —mediana en
  vez de umbral fijo porque la escala es propia de cada activo \\
Drawdown & deteriorado / sano & \(\text{drawdown\_historico} \leq -0.20\) \(\Rightarrow\)
  deteriorado \\
\bottomrule
\end{tabularx}
\caption{Las 3 dimensiones de régimen de Q2-13.}
\end{table}

Clasificación (misma regla que Q2-9/Q2-10): \emph{cambio de signo} si cambia entre
categorías; \emph{dependiente de régimen} si no cambia de signo pero la dispersión entre
categorías es alta; \emph{consistente} en otro caso. Una relación \emph{consistente} entre
regímenes es más confiable que una con Pearson global más alto pero que cambia de signo
—el global puede promediar dos relaciones opuestas en distintos regímenes—.

\subsection{Síntesis descriptiva y nivel de evidencia (Q2-11)}

Consolida, por variable/target, todo lo anterior en un \textbf{nivel de evidencia
descriptiva} explícito —regla, no un score de ML—:

\begin{table}[h]
\centering
\small
\begin{tabularx}{\linewidth}{@{}l X@{}}
\toprule
\textbf{Nivel} & \textbf{Condición} \\
\midrule
\code{alta} & \(|\text{pearson}| \geq 0.10\) (fuerza al menos \dq{débil}) \textbf{y}
  estable en el tiempo \textbf{y} consistente entre los 3 regímenes \\
\code{media} & tiene fuerza, pero falla en estabilidad temporal \textbf{o} en régimen (no
  en ambas) \\
\code{baja} & sin fuerza, o con fuerza pero inestable en el tiempo \textbf{y} dependiente
  de régimen a la vez \\
\bottomrule
\end{tabularx}
\caption{Regla explícita de \code{nivel\_evidencia\_descriptiva} (Q2-11).}
\end{table}

\begin{advertencia}
Esta tabla \textbf{no es un ranking predictivo}: evidencia descriptiva sólida no garantiza
poder predictivo fuera de muestra — eso se valida con backtesting en una etapa posterior
(Sección~\ref{ssec:proximos-pasos}).
\end{advertencia}

\subsubsection*{Lectura de un resultado real: BTCUSD vs. ETHUSD}

La corrida \code{Q2-EXEC} (validación de punta a punta, 2026-10-12) extendió
\code{valores} a los 6 activos operados por primera vez —hasta entonces cada ítem se había
validado solo contra BTCUSD—, sobre datos frescos de Windows (a diferencia del snapshot
desactualizado de Mac usado para ilustrar las limitaciones de \(x2\_score\) en la
Sección~\ref{sec:q1}).

\begin{table}[h]
\centering
\small
\begin{tabular}{@{}l r r r@{}}
\toprule
Activo & \code{alta} & \code{media} & \code{baja} \\
\midrule
BTCUSD & 26 & 5 & 95 \\
ETHUSD & 25 & 9 & 92 \\
\bottomrule
\end{tabular}
\caption{Distribución de \code{nivel\_evidencia\_descriptiva} contra
\code{retorno\_log\_1h}, de las 126 variables analizadas en cada activo.}
\end{table}

Dos lecturas de ese resultado: primero, la distribución es \textbf{razonable y similar}
entre dos activos de la misma familia (ambos \(\approx\)20\% \code{alta}, resto
mayormente \code{baja}) — señal de que el pipeline generaliza entre activos en vez de
producir ruido activo-específico. Segundo, el top-8 de variables \code{alta} contra
\code{retorno\_log\_1h} coincide \emph{casi en el mismo orden} entre ambos activos:
\code{x2\_score\_cross}, \code{fundamental\_estabilidad\_90d}, \code{x2\_score},
\code{x3\_ema\_dist\_12}, \code{distancia\_ema\_12}, \code{ruptura\_rango},
\code{x3\_bb\_pos}, \code{x3\_sma\_dist\_20}.

\begin{nota}
Que \(x2\_score\) y sus derivadas aparezcan acá con evidencia \code{alta} —a diferencia de
la limitación de \(x2\_score\) 100\% \code{NaN} documentada en la
Sección~\ref{sec:q1}— no es una contradicción: esa limitación es específica del snapshot
local de \code{Data/}/\code{x2\_history.json} en Mac (desfasado en fechas); la corrida
\code{Q2-EXEC} se hizo en Windows con datos frescos y solapados, donde \(x2\_score\) sí
varía.
\end{nota}

\subsection{KPIs adicionales: información mutua y correlación de distancia (Q2-14)}

Objetivo explícito: evaluar \textbf{caso a caso} —no adoptar por default— si una métrica
adicional aporta algo que Pearson/Spearman/coseno no muestren. Solo dos candidatas, ninguna
más (principio de parsimonia).

\begin{itemize}[leftmargin=1.4em,itemsep=2pt]
\item \textbf{Información mutua} (\code{sklearn.mutual\_info\_regression},
  \code{n\_neighbors=3}): capta cualquier dependencia, lineal o no; se reporta en nats
  (no acotada en \([0,1]\)) \(\Rightarrow\) comparación por \textbf{ranking}, no valor
  absoluto —\code{entra\_top\_mi\_sin\_estar\_en\_pearson} si entra al top-10 de
  información mutua sin estar en el top-10 de \(|\text{Pearson}|\)—.
\item \textbf{Correlación de distancia} (Székely–Rizzo), acotada en \([0,1]\)
  \(\Rightarrow\) comparable directo con \(|\text{Pearson}|\):
\end{itemize}

\begin{align}
A_{ij} &= a_{ij} - \bar{a}_{i\cdot} - \bar{a}_{\cdot j} + \bar{a}
  \qquad (\text{doble centrado}),\quad a_{ij} = |x_i - x_j| \\
\mathrm{dCov}^2 &= \overline{A \cdot B}, \qquad \mathrm{dVar}_X^2 = \overline{A\cdot A},
  \qquad \mathrm{dVar}_Y^2 = \overline{B \cdot B} \\
\mathrm{dcor} &= \sqrt{\frac{\mathrm{dCov}^2}{\sqrt{\mathrm{dVar}_X^2 \cdot
  \mathrm{dVar}_Y^2}}}
\end{align}

\(O(n^2)\) en tiempo/memoria \(\Rightarrow\) calculada sobre una \textbf{submuestra fija}
de 3.000 filas (semilla 42), nunca sobre la tabla completa. Se marca
\code{aporta\_valor\_dcor} si \(\mathrm{dcor} - |\mathrm{Pearson}| > 0.10\).

\bigsep
\textbf{Resultado real en BTCUSD}: varias variables de volatilidad/liquidez
(\code{x3\_vol\_24h}, \code{x3\_atr\_pct}, \code{x3\_bb\_width}, \code{tick\_volume\_z},
\code{spread\_relativo}, \code{volumen\_relativo}) muestran \(\mathrm{dcor}\)
\(0.17\)–\(0.34\) contra \code{retorno\_log\_1h}, muy por sobre su Pearson casi nulo —
dependencia no lineal real, aunque sin definir signo ni servir de validación. Este CSV
queda deliberadamente fuera de la síntesis (Q2-11) y del export principal (Q2-5) — anexo
exploratorio, no insumo obligatorio.

\subsection{Export estructurado de métricas (Q2-5)}

No recalcula nada: recopila en un único JSON anidado \code{\{variable: \{...\}\}} todo lo
ya calculado y guardado en CSV por las subsecciones anteriores. Un archivo por activo:
\code{\{valor\}\_q2\_5\_export\_metricas.json} — el \textbf{único archivo} que
\code{X5\_Q3.ipynb} lee de Q2, pensado para alimentar después un análisis por IA que
concluya cosas como \dq{el drawdown se relaciona con \code{N}} (texto literal del
notebook).

\begin{lstlisting}
{
  "drawdown_7d": {
    "targets": {"Close": {...}, "retorno_log_1h": {
        "pearson": ..., "spearman": ..., "coseno": ..., "direccion_fuerza": ...,
        "no_linealidad": {"r2_lin": ..., "r2_cuad": ..., "delta_r2": ..., "forma": ...},
        "estabilidad": {"clasificacion": ..., "pearson_std_entre_anios": ...},
        "nivel_evidencia_descriptiva": ..., "cobertura_pct": ...
    }},
    "ceteris_paribus": {...}, "retorno_pasado": {...},
    "regimen": {...}, "redundancia": {...}
  }
}
\end{lstlisting}

\newpage
% ═══════════════════════════════════════════════════════════════════════════
\section{X5\_Q3 — hipótesis parámetro$\leftrightarrow$variable y tabla final}
\label{sec:q3}
% ═══════════════════════════════════════════════════════════════════════════

Tercer y último notebook. Consume \code{\{valor\}\_tabla\_estado\_v2.csv} (Q1) y
\code{\{valor\}\_q2\_5\_export\_metricas.json} (Q2) — gate de arranque: requiere Q1 y Q2 ya
completados.

\subsection{Cómo correr el notebook}

\begin{enumerate}[leftmargin=1.4em,itemsep=2pt]
\item Verificar primero que existan \code{\{valor\}\_tabla\_estado\_v2.csv} (Q1) y
  \code{\{valor\}\_q2\_5\_export\_metricas.json} (Q2) para cada activo de \code{valores}
  —ambos levantan \code{FileNotFoundError} explícito si faltan—.
\item Es el notebook más corto de los tres (17 celdas): la mayoría son celdas markdown
  con la hipótesis de cada parámetro (sin código, documentación pura) — solo la tabla
  final (\code{Q3-4}) y la vista unificada (\code{Q3-8}) ejecutan código.
\item Ejecutar de arriba hacia abajo. La tabla final (\code{Q3-4}) consulta
  programáticamente \code{export\_metricas} para cada variable candidata ya fijada en las
  celdas markdown anteriores — no recalcula ninguna métrica de Q2.
\item Verificar al final \code{\{valor\}\_q3\_4\_tabla\_final.csv} y
  \code{\{valor\}\_q3\_8\_vista\_unificada.csv}, y que el \code{.zip} los incluye a ambos.
\end{enumerate}

\subsection{Alcance: lo que Q3 sí y no hace}

\begin{table}[h]
\centering
\small
\begin{tabularx}{\linewidth}{@{}X X@{}}
\toprule
\textbf{Sí hace} & \textbf{No hace} \\
\midrule
Produce hipótesis \code{parámetro = f(variable(s))} apoyadas en el pre-registro y la
evidencia de Q2; una vista unificada de variables relevantes + precio; una tabla final de
tuplas con evidencia y prioridad de prueba. &
No implementa nada en \code{X0\_data\_supports.py}/\code{X1\_trading.py} de producción. No
valida con backtesting — etapa posterior sobre \code{X4\_backtester.py}. \\
\bottomrule
\end{tabularx}
\end{table}

\begin{clave}
\dq{Una hipótesis con evidencia fuerte en Q2 sigue siendo una hipótesis hasta que el
backtesting la confirme.} —cita literal del notebook.
\end{clave}

\subsection{Por qué K, N\_EXP y LAMBDA requieren dos niveles de análisis}

Q2 corre sobre precio/variables de estado, nunca sobre soportes ni órdenes — aporta solo
evidencia indirecta sobre por qué convendría variar cada parámetro de X0:

\begin{enumerate}[leftmargin=1.4em,itemsep=2pt]
\item \textbf{Nivel 1 — efecto sobre los soportes (X0)}: cantidad, calidad individual,
  recencia, cobertura del rango, concentración, distancia entre niveles, simetría de
  aislamiento, estabilidad.
\item \textbf{Nivel 2 — efecto posterior en X1}: órdenes creadas/activadas, cierres,
  dinero cerrado, MAE/MFE (\emph{Maximum Adverse/Favorable Excursion}), duración,
  estabilidad por período.
\end{enumerate}

\begin{advertencia}
\dq{No debe concluirse que un valor es bueno solo porque genera soportes visualmente bien
distribuidos} —debe demostrar utilidad operativa en X1, no solo verse bien en X0.
\end{advertencia}

\subsection{Mecánica exacta de cada parámetro}

\begin{table}[h]
\centering
\small
\begin{tabularx}{\linewidth}{@{}l l X@{}}
\toprule
\textbf{Parámetro} & \textbf{Dónde actúa} & \textbf{Fórmula / mecánica} \\
\midrule
\code{N} & X1 & cantidad de soportes activos mantenidos por X1
  (\code{n\_sizes\_ejecucion}) \\
\code{A} & X1 & doble rol: filtro de OE \((P_0-P_i)\cdot L \geq A\) (lotaje
  \textbf{mínimo}) y umbral de ganancia para el primer SL del trailing
  \((P_0-P_i)\cdot L \geq A\) (lotaje \textbf{real} de la posición) \\
\code{factor\_B} & X1 & \(B = \text{factor\_B} \cdot A\) — distancia del trailing una
  vez activo, \(\mathrm{sl}_{\mathrm{nuevo}} = P_0 - B/L\) \\
\code{K} & X0 & \(y = \text{Low}_{\mathrm{left}} + \text{High}_{\mathrm{left}} +
  K\cdot(\text{Low}_{\mathrm{right}} + \text{High}_{\mathrm{right}})\) — peso del
  aislamiento futuro vs. pasado \\
\code{N\_EXP} & X0 & \(w = t^{N\_EXP}\), \(t\in[0,1]\) — exponente de recencia \\
\code{LAMBDA} & X0 & \(\mathrm{FO} = \overline{z} - \mathrm{LAMBDA}\cdot
  \mathrm{cv}(H_n)\) — penalización por dispersión desigual entre soportes \\
\bottomrule
\end{tabularx}
\caption{Mecánica de cada parámetro Y, según \code{docs/context/pars\_configurables.md}.}
\end{table}

\begin{nota}
\textbf{Dualidad de roles de \code{A}} (por qué su prioridad de prueba es \dq{media}, no
\dq{alta}): el filtro de OE y la activación del trailing tiran en direcciones distintas
frente a un mismo contexto —con drawdown profundo, un \code{A} alto reduce entradas nuevas
pero también retrasa la protección de una posición ganadora—. El efecto neto del drawdown
sobre \code{A} no tiene dirección única; podría requerir separar
\code{A\_OE}/\code{A\_TRAILING} más adelante (decisión pendiente).
\end{nota}

\subsubsection*{Derivación: el umbral de \code{A} es la misma fórmula en los dos roles}

Vale la pena derivar esto con álgebra —ya resuelto en
\code{docs/plans/formulas.tex}, resumido acá porque explica por qué \code{A} es un único
parámetro con dos efectos, no dos parámetros independientes—. Sea \(m\) el lote mínimo del
bróker, \(U\) las unidades por lote, \(M\) el multiplicador de lote configurado
(\code{LOTAJES\_M}).

\textbf{Caso 1 — O0 se convierte en OE} (\code{crear\_ordenes\_espera}): se crea el buy
limit cuando \((P_0-P_i)\,m\,U \geq A\) —con el lotaje \textbf{mínimo} del activo, no el
configurado—:

\begin{align}
P_0 \;\geq\; P_i + \frac{A}{mU}
\end{align}

\textbf{Caso 2 — primera activación del Stop Loss} (\code{trailing\_stop}): opera sobre el
lotaje \textbf{real} de la posición (\code{orden.volume} \(=Mm\)); el umbral efectivo
\(a_{\mathrm{efectivo}} = A\cdot M\) y la ganancia \((P_0-P_i)\,Mm\,U\) comparten el mismo
factor \(M\), que se cancela al dividir:

\begin{align}
(P_0-P_i)\,MmU \;\geq\; AM \quad\Longrightarrow\quad P_0 \;\geq\; P_i + \frac{A}{mU}
\end{align}

\begin{clave}
\textbf{Mismo umbral de precio en los dos roles} —\(P_i + A/(mU)\)—, aunque \(P_i\)
signifique algo distinto en cada caso (el soporte en el Caso 1, el precio de apertura de
la OA en el Caso 2). \(M\) (\code{LOTAJES\_M}) \textbf{no aparece} en ninguno de los dos:
se cancela en el Caso 2, y se excluyó deliberadamente del Caso 1 el 2026-09-25 (antes
usaba el lotaje configurado, no el mínimo — subir \(M\) relajaba el gap de activación sin
motivo). Esta es la raíz algebraica de por qué separar \code{A\_OE}/\code{A\_TRAILING}
—la decisión pendiente señalada arriba— cambiaría una fórmula hoy unificada en dos.
\end{clave}

\begin{table}[h]
\centering
\small
\begin{tabularx}{\linewidth}{@{}l r r r X@{}}
\toprule
\textbf{Activo} & \(m\) & \(U\) & \(A/(mU)\) & \textbf{Gap de activación resultante} \\
\midrule
BTCUSD & 0{,}01 & 1   & 300{,}00 & soporte en 65.000 \(\to\) OE recién a 65.300 \\
ETHUSD & 0{,}10 & 1   & 30{,}00  & soporte en 3.000 \(\to\) OE recién a 3.030 \\
TSLA   & 0{,}01 & 100 & 3{,}00   & soporte en 250 \(\to\) OE recién a 253 \\
GOOGL/NVDA/AMZN & 0{,}01 & 100 & 3{,}00 & mismo gap que TSLA (idéntico \(mU\)) \\
\bottomrule
\end{tabularx}
\caption{Gap de activación \(A/(mU)\) con \(A=3\) (valor vigente, igual para los 6
activos). Valores de \(P_i\) ilustrativos, tomados de \code{docs/plans/formulas.tex}.}
\end{table}

\textbf{Lectura}: el mismo \(A=3\) produce un gap de \$300 en BTCUSD pero de solo \$3 en
las acciones —porque \(U=100\) en acciones multiplica por 100 cada dólar de movimiento de
precio, mientras que en cripto \(U=1\)—. Cualquier hipótesis \(A=f(\text{volatilidad})\)
de la Sección~\ref{sec:q3} debe leerse sobre esta escala ya normalizada por \(mU\), no
sobre el valor crudo de \(A\) en USD.

\subsection{Tabla final de hipótesis}

\begin{table}[h]
\centering
\small
\begin{tabularx}{\linewidth}{@{}l X X l@{}}
\toprule
\textbf{Parámetro} & \textbf{Variables candidatas} & \textbf{Relación hipotética} &
  \textbf{Prioridad} \\
\midrule
\code{N} & \code{x3\_drawdown\_20}, \code{x3\_drawdown\_50}, \code{drawdown\_7d} &
  drawdown sano (menos profundo) \(\to\) \code{N} alto; drawdown profundo \(\to\)
  \code{N} bajo & \textbf{alta} \\
\code{A} & \code{x3\_atr\_pct}, \code{x3\_bb\_width} &
  volatilidad alta \(\to\) \code{A} alto; volatilidad baja \(\to\) \code{A} bajo
  (dirección neta sobre drawdown ambigua) & media \\
\code{factor\_B} & \code{x3\_atr\_pct}, \code{adx\_14}, \code{regimen\_tendencia} &
  volatilidad alta o tendencia fuerte \(\to\) \code{factor\_B} alto; volatilidad baja
  \(\to\) \code{factor\_B} bajo & media \\
\code{K} & \code{duracion\_en\_rango}, \code{ruptura\_rango} &
  niveles persistentes/mercado estable \(\to\) \code{K} alto & baja \\
\code{N\_EXP} & \code{cambio\_volatilidad}, \code{frecuencia\_extremos\_30d} &
  cambio de régimen o volatilidad creciente \(\to\) \code{N\_EXP} alto & media \\
\code{LAMBDA} & \code{x3\_bb\_width}, \code{duracion\_en\_rango} (junto con \code{N}) &
  rango amplio/lateral con soportes concentrados \(\to\) \code{LAMBDA} alto &
  baja (bloqueada) \\
\code{LOTAJES\_M} & — & sin hipótesis dinámica — fijo (rango degenerado \((1,1)\)) & — \\
\code{PERDIDA\_MAX} & — & pausado — en evaluación para una capa externa de gestión de
  riesgo & — \\
\bottomrule
\end{tabularx}
\caption{Tabla final de tuplas (parámetro, variable candidata, relación hipotética,
prioridad de prueba) — \code{Q3-4}.}
\end{table}

\textbf{Criterio de selección de variables}: por parámetro, la de mejor evidencia
descriptiva disponible (Q2-11/Q2-14) que no fuera redundante con otra ya elegida para el
mismo parámetro (Q2-12) — ej. \code{x3\_vol\_24h} se descarta en \code{A} por ser
redundante con \code{x3\_atr\_pct} (Pearson \(0.91\)). Ninguna tupla está validada — son el
insumo de entrada, no una conclusión.

\begin{advertencia}
\textbf{Bloqueador de \code{LAMBDA}, sin resolver}: el rango de exploración de X5
documentado (\(0.001\)–\(0.02\)) \textbf{no cubre} el valor productivo real (\(0.2\)) —
no deben sacarse conclusiones de una exploración que no contiene el valor de referencia.
\end{advertencia}

\textbf{Prioridad de prueba} es juicio explícito, no un score automático de la evidencia
Q2: \code{A}/\code{factor\_B}/\code{K}/\code{N\_EXP}/\code{LAMBDA} no son hipótesis de
predicción de retorno —su evidencia Q2 baja contra \code{retorno\_log\_1h}
(\(|\text{Pearson}|<0.02\) para \code{x3\_atr\_pct}/\code{x3\_bb\_width}/\code{adx\_14})
es \textbf{esperada}, no indica menor madurez—. Se basa en: madurez de la fórmula
hipotética, bloqueadores documentados y el rol de cada parámetro en la mecánica de X0/X1.

\subsection{Vista unificada (Q3-8)}

Tabla con \code{DateTime}, \code{Close} y las variables cuyo
\code{nivel\_evidencia\_descriptiva} es \code{alta} o \code{media} contra
\code{retorno\_log\_1h} (el target principal — \code{Close} es solo referencial).
Guardada en \code{\{valor\}\_q3\_8\_vista\_unificada.csv}, pensada para graficar/revisar a
ojo, \textbf{no} para modelar.

\subsection{Próximos pasos}
\label{ssec:proximos-pasos}

Con \code{Q3-EXEC} satisfecho la implementación de los tres notebooks queda cerrada — lo
que sigue vive \textbf{fuera} de X5\_alt\_v2:

\begin{enumerate}[leftmargin=1.4em,itemsep=3pt]
\item \textbf{Puente a backtesting real sobre \code{X4\_backtester.py}.} Ninguna hipótesis
  de la Tabla anterior está confirmada — son candidatas a probar, no reglas.
\item \textbf{Orden de validación sugerido}: empezar por \code{N} (prioridad
  \code{alta}, hipótesis madura, sin bloqueadores) antes que \code{A}/\code{factor\_B}/
  \code{N\_EXP} (prioridad \code{media}) y \code{K}/\code{LAMBDA} (prioridad
  \code{baja}/bloqueada).
\item \textbf{Decisiones de diseño pendientes antes de testear} (no son trabajo de Q3):
  separar \code{A\_OE}/\code{A\_TRAILING}, y decidir si el rango de exploración de
  \code{LAMBDA} en \code{config.py} está desactualizado.
\item \textbf{No expandir a otros activos} (\code{ETHUSD}, \code{TSLA}, \code{GOOGL},
  \code{NVDA}, \code{AMZN}) hasta validar el enfoque con backtesting sobre BTCUSD.
\end{enumerate}

\begin{nota}
El TO DO más cercano a esta etapa es \dq{X5 — estudio posterior de parámetros manipulables
y variables explicativas} (\code{docs/tracking/todos.md}, sección X5\_alt) — hoy escrito en
términos de la versión anterior, antes de que existiera \code{X5\_Q3}. Pide exactamente la
misma tabla que \code{Q3-4} ya produjo — debería actualizarse para apuntar a
\code{\{valor\}\_q3\_4\_tabla\_final.csv} en vez de pedir reconstruirla desde cero.
\end{nota}

\newpage
% ═══════════════════════════════════════════════════════════════════════════
\section{Preguntas frecuentes}
\label{sec:faq}
% ═══════════════════════════════════════════════════════════════════════════

\paragraph{¿Por qué no basta con mirar Pearson contra \code{Close}?}
Porque \code{Close} no es estacionario —BTCUSD recorrió \(\sim\)20.000 a \(\sim\)100.000
USD entre 2021 y 2026—, así que casi cualquier variable con tendencia propia va a
\dq{correlacionar} con el nivel de precio solo por coincidir en el tiempo, no por aportar
información real. Por eso todo el núcleo descriptivo (Sección~\ref{sec:q2}) usa el retorno
contemporáneo (\(\text{retorno\_log\_1h}\)) como target principal, y el nivel como
referencia secundaria (Sección~\ref{ssec:dos-paneles}).

\paragraph{¿El coseno no es redundante con Pearson, si los dos miden \dq{asociación}?}
No: Pearson centra ambas series en su propia media antes de medir covariación
(\(\mathrm{corr}(x,y)=\cos(x-\bar{x},y-\bar{y})\)); el coseno de la
Sección~\ref{ssec:coseno} no centra nada. Contra el nivel de precio (Panel A) esa
diferencia puede incluso invertir el signo (Tabla~\ref{tab:coseno-ejemplo}); contra el
retorno (Panel B) coincide en signo pero se atenúa cuando la variable tiene un baseline
lejos de cero (ej. \code{x3\_rsi\_14}, centrado en \(\sim\)50). Son dos preguntas
distintas, no la misma pregunta dos veces.

\paragraph{¿Por qué \code{K}, \code{N\_EXP} y \code{LAMBDA} necesitan backtesting en dos
niveles y no solo la evidencia de Q2?}
Porque actúan \textbf{dentro del algoritmo de soportes de X0} —generan qué soportes
existen, no gestionan órdenes directamente—, y su efecto sobre el retorno es indirecto y
desfasado en el tiempo: un soporte calculado hoy puede producir una orden que cierra
varios días después. Q2 nunca corre sobre soportes, solo sobre precio/variables de
estado —aporta evidencia indirecta de por qué convendría variar el parámetro, no de si
realmente mejora el resultado en X0 o en X1 (Sección~\ref{sec:q3}, nivel 1 vs. nivel 2).

\paragraph{¿Una variable con \code{nivel\_evidencia\_descriptiva = alta} ya está lista
para usarse en producción?}
No. \dq{Alta} significa evidencia descriptiva sólida —fuerza, estabilidad temporal y
consistencia entre regímenes (Sección~\ref{sec:q2})—, no poder predictivo fuera de
muestra ni validación con backtesting. La frase que cierra \code{X5\_Q3.ipynb} lo dice sin
ambigüedad: \dq{una hipótesis con evidencia fuerte en Q2 sigue siendo una hipótesis hasta
que el backtesting la confirme.}

\paragraph{¿Por qué \code{A} y \code{factor\_B} tienen evidencia Q2 casi nula contra el
retorno, si son justamente los parámetros que se quiere hacer dinámicos?}
Porque no son hipótesis de \emph{predicción de retorno} sino de \emph{escala de la
mecánica de órdenes} frente al estado del activo —\(A\) regula a qué distancia se activa
una orden o un trailing, no si el precio va a subir o bajar—. Un Pearson casi nulo contra
\code{retorno\_log\_1h} es exactamente lo esperado para este tipo de parámetro
(Sección~\ref{sec:q3}), y no implica que la hipótesis esté peor fundamentada que la de
\code{N}.

\paragraph{¿Por qué no se expande a los otros 5 activos todavía?}
Porque el enfoque —tabla de estado, núcleo descriptivo, hipótesis— se valida primero sobre
BTCUSD antes de generalizar, mismo criterio que ya regía para la versión anterior. La
corrida \code{Q2-EXEC} sí probó los 6 activos puntualmente (Sección~\ref{sec:q2}) para
verificar que el pipeline generaliza, pero eso es distinto de adoptar los 6 como alcance
corriente de los tres notebooks.

\paragraph{¿Qué significa que una variable se \dq{omita con aviso}?}
Que para ese par variable/target no se alcanzó el mínimo de 30 observaciones no nulas, o
que una de las dos series no tiene dispersión (es constante) — el mismo filtro en las
catorce subsecciones de \code{X5\_Q2.ipynb} (Sección~\ref{sec:q2}). En el snapshot de Mac
usado para desarrollar estos notebooks, esto ocurre sobre todo con \(x2\_score\) y sus
derivadas —100\% \code{NaN} para BTCUSD, por el desfase de fechas ya documentado en la
Sección~\ref{sec:q1}—. No es un error del código: es el notebook reportando honestamente
que, con los datos disponibles en ese momento, no hay evidencia suficiente para esa
variable — se corrige solo con datos más frescos, no con un cambio de lógica.

\newpage
% ═══════════════════════════════════════════════════════════════════════════
\section{Síntesis}
% ═══════════════════════════════════════════════════════════════════════════

\begin{itemize}[leftmargin=1.4em,itemsep=4pt]
\item \code{X5\_Q1.ipynb} construye una tabla H1 continua del \emph{estado del activo} —
  drawdown, volatilidad, ruido, tendencia/régimen, liquidez, fundamental— reutilizando los
  indicadores de X3 en vez de recalcularlos, y documentando cada variable en un catálogo
  auditable.
\item \code{X5\_Q2.ipynb} cuantifica, con evidencia estadística explícita (Pearson,
  Spearman, coseno, no linealidad, ceteris paribus, estabilidad, régimen, redundancia), qué
  tan confiable es cada relación variable$\leftrightarrow$precio/retorno — y resume todo en un nivel de
  evidencia descriptiva por variable.
\item \code{X5\_Q3.ipynb} traduce esa evidencia en hipótesis explícitas
  \(\text{parámetro} = f(\text{variable}(s))\) para \code{N}, \code{A}, \code{factor\_B},
  \code{K}, \code{N\_EXP} y \code{LAMBDA}, con prioridad de prueba justificada —
  \textbf{sin implementar ninguna como regla}.
\item El puente hacia backtesting real sobre \code{X4\_backtester.py} es la etapa
  siguiente, explícitamente fuera del alcance de los tres notebooks documentados aquí.
\end{itemize}

\subsection*{La cadena de dependencias, de punta a punta}

Los tres notebooks forman una cadena estricta —cada uno solo puede correr después de que
el anterior haya terminado de punta a punta, nunca en paralelo—:

\begin{table}[h]
\centering
\small
\begin{tabularx}{\linewidth}{@{}l X X l@{}}
\toprule
\textbf{Notebook} & \textbf{Consume} & \textbf{Produce} & \textbf{Gate} \\
\midrule
\code{X5\_Q1} & \code{Data/}, \code{resources/x2/}, \code{resources/x3/} (vía X3) &
  \code{\{valor\}\_tabla\_estado\_v2.csv}, \code{\{valor\}\_catalogo\_variables.csv} &
  \code{Q1-EXEC} \\
\code{X5\_Q2} & \code{\{valor\}\_tabla\_estado\_v2.csv} &
  \textasciitilde19 CSV por ítem, \code{\{valor\}\_q2\_5\_export\_metricas.json} &
  \code{Q2-EXEC} \\
\code{X5\_Q3} & \code{\{valor\}\_tabla\_estado\_v2.csv},
  \code{\{valor\}\_q2\_5\_export\_metricas.json} &
  \code{\{valor\}\_q3\_8\_vista\_unificada.csv}, \code{\{valor\}\_q3\_4\_tabla\_final.csv} &
  \code{Q3-EXEC} \\
\bottomrule
\end{tabularx}
\caption{Insumos, productos y gate de ejecución de cada notebook — la referencia rápida
para saber qué correr antes de qué.}
\end{table}

Con los tres gates satisfechos (estado a 2026-10-12, Sección~\ref{ssec:proximos-pasos}),
la implementación de \code{X5\_alt\_v2} queda cerrada: lo único pendiente es la
validación con backtesting, fuera del alcance de estos tres notebooks.

\appendix
\section*{Anexo: catálogo completo de fórmulas (X5\_Q1)}
\addcontentsline{toc}{section}{Anexo: catálogo completo de fórmulas (X5\_Q1)}

Referencia consolidada de todas las variables que \code{X5\_Q1.ipynb} agrega sobre la
tabla base de precio+X3+X2 (Sección~\ref{sec:q1}) — complementa, no reemplaza, el catálogo
en vivo que el propio notebook genera por activo
(\code{\{valor\}\_catalogo\_variables.csv}, con cobertura/nulos/fechas reales).

\begin{longtable}{@{}p{2.3cm} p{3.4cm} p{7.0cm} p{2.0cm}@{}}
\caption{Catálogo completo de fórmulas agregadas por \code{X5\_Q1.ipynb}, por familia.}
\\
\toprule
\textbf{Familia} & \textbf{Variable} & \textbf{Fórmula} & \textbf{Ventana} \\
\midrule
\endfirsthead
\toprule
\textbf{Familia} & \textbf{Variable} & \textbf{Fórmula} & \textbf{Ventana} \\
\midrule
\endhead
\bottomrule
\endfoot
Técnicos transf. & \code{distancia\_sma}\(_w\) / \code{distancia\_ema}\(_w\) &
  \((\text{Close}-\text{media}_w)/\text{ATR}_{14}\) & \(w{=}\)20/50/200 (sma),
  12/26 (ema) \\
Técnicos transf. & \code{pendiente\_sma}\(_w\) / \code{pendiente\_ema}\(_w\) &
  \((\text{media}_w.\mathrm{diff}(w)/w)/\text{ATR}_{14}\) & misma \(w\) \\
Drawdown & \code{drawdown\_historico} &
  \((\text{Close}-\mathrm{cummax})/\mathrm{cummax}\) & histórico \\
Drawdown & \code{duracion\_drawdown} & velas desde el último máximo histórico & — \\
Drawdown & \code{recuperacion\_desde\_minimo} &
  \((\text{Close}-\min_{\mathrm{ep}})/\min_{\mathrm{ep}}\) & episodio actual \\
Drawdown & \code{velocidad\_caida} &
  \(\text{drawdown\_historico}/\text{duracion\_drawdown}\) & — \\
Drawdown & \code{drawdown}\(_d\), \(d{\in}\{7,30,90,365\}\)d &
  \((\text{Close}-\max_h)/\max_h\) & 168/720/2160/8760h \\
Volatilidad & \code{cambio\_volatilidad} & \(x3\_vol\_24h / x3\_vol\_7d\) & — \\
Volatilidad & \code{vol\_alcista}\(_{30d}\) / \code{vol\_bajista}\(_{30d}\) &
  \(\mathrm{std}_{30d}(r_t{>}0\text{ o }{<}0)\cdot\sqrt{8760}\) & 720h \\
Volatilidad & \code{frecuencia\_extremos}\(_{30d}\) &
  \(\mathrm{media}_{30d}(|r_t|{>}2\sigma_{\mathrm{exp}})\) & 720h,
  \(\sigma\) \texttt{expanding} \\
Liquidez & \code{spread\_relativo} &
  \(\text{Spread}/\mathrm{media}_{30d}(\text{Spread})\) & 720h \\
Liquidez & \code{tick\_volume\_z} &
  \((\text{Tick\_Volume}-\mathrm{media}_{30d})/\mathrm{std}_{30d}\) & 720h \\
Liquidez & \code{volumen\_relativo} (RVOL) &
  \(\text{Tick\_Volume}/\mathrm{media}_{30d}(\text{Tick\_Volume})\) & 720h \\
Fundamental & \code{fundamental\_cambio} & \(x2\_score.\mathrm{diff}()\) & — \\
Fundamental & \code{fundamental\_direccion\_cambio} &
  \(\mathrm{sign}(\text{fundamental\_cambio})\) & — \\
Fundamental & \code{fundamental\_magnitud\_cambio} &
  \(|\text{fundamental\_cambio}|\) & — \\
Fundamental & \code{fundamental\_edad} & velas desde la última actualización real & — \\
Fundamental & \code{fundamental\_frecuencia\_efectiva}\(_{90d}\) &
  \(\mathrm{media}_{90d}(\text{actualización real})\) & 2160h \\
Fundamental & \code{fundamental\_estabilidad}\(_{90d}\) & \(\mathrm{std}_{90d}(x2\_score)\)
  & 2160h \\
Fundamental & \code{fundamental\_contradiccion} &
  \((x2\_score\_cross{-}0{,}5)(x2\_score\_tendencia{-}0{,}5)\) & — \\
Ruido & \code{ruido\_cambios\_signo}\(_{30d}\) &
  \(\mathrm{media}_{30d}(\mathrm{sign}(r_t)\neq\mathrm{sign}(r_{t-1}))\) & 720h \\
Ruido & \code{ruido\_ratio\_movimiento}\(_{30d}\) &
  \(\sum|r_t| / |\sum r_t|\) (inverso Kaufman) & 720h \\
Ruido & \code{ruido\_mecha}\(_{30d}\) &
  \(\mathrm{media}((\text{H-L-}|\text{C-O}|)/\text{H-L})\) & 720h \\
Ruido & \code{ruido\_atr\_sin\_direccion}\(_{30d}\) &
  \(\mathrm{media}(x3\_atr\_pct)/|\sum r_t|\) & 720h \\
Tendencia/\\régimen & \code{adx\_14} & \(\mathrm{WilderEMA}(\mathrm{DX},14)\) & 14 \\
Tendencia/\\régimen & \code{regimen\_tendencia} &
  \code{sideways} si \(\mathrm{ADX}{<}20\); si no, signo de
  \(x3\_trend\_slope\_50\) & — \\
Tendencia/\\régimen & \code{ruptura\_rango} &
  \(+1/-1/0\) vs. máx./mín. de 20 velas previas & 20 \\
Tendencia/\\régimen & \code{duracion\_en\_rango} & velas desde la última ruptura & — \\
\end{longtable}

\section*{Anexo: mapa de archivos}
\addcontentsline{toc}{section}{Anexo: mapa de archivos}

\begin{lstlisting}
scripts/
  X5_Q1.ipynb          # tabla de estado + catalogo
  X5_Q2.ipynb          # core descriptivo (14 items)
  X5_Q3.ipynb          # hipotesis + tabla final
  x5_alt_export.py     # exportar_zip(nombre_notebook, archivos)
  otros_strings/
    X5_P1.ipynb        # v1, congelado
    X5_P2.ipynb        # v1, congelado

resources/x5_alt/      # TODO fuera de git -- generado en Windows
  {valor}_tabla_estado_v2.csv
  {valor}_catalogo_variables.csv
  {valor}_q2_*.csv              # 1 archivo por item Q2-3..Q2-14
  {valor}_q2_5_export_metricas.json
  {valor}_q3_8_vista_unificada.csv
  {valor}_q3_4_tabla_final.csv

docs/plans/
  X5_alternativo.md              # contexto y estado de avance, seccion 10
  x5_contexto_conversacion.md    # documento rector
docs/context/
  pars_configurables.md          # definicion exacta de cada parametro Y
docs/tracking/
  todos.md   done.md             # seccion X5_alt_v2
\end{lstlisting}

\end{document}
